\section{Discussion}

\subsection{Interpreting the Korea--Japan Divergence}

The central empirical finding is a sharp asymmetry in the valuation consequences of Japan's three governance reform milestones. In the primary MSCI EM Asia--TOPIX specification, Japan's two early soft-law interventions---the 2014 Stewardship Code and the 2015 Corporate Governance Code---produced positive abnormal spread movements, indicating that the entire EM Asia regional basket held its position against Japan in the immediate post-reform periods. The 2023 TSE P/B Reform Request is categorically different: the MSCI EM Asia--TOPIX spread declined monotonically to $-2.584$ P/B points by $\tau = 20$, and the secondary KOSPI--TOPIX spread widened to $-6.481$ P/B points by $\tau = 24$---the largest event in the sample. Japan's mandatory P/B reform lifted TOPIX above the entire regional peer basket, with Korea's individual shortfall exceeding the regional average.

This asymmetry carries a specific interpretation. Soft governance codes---whether Japan's 2014 Stewardship Code or its 2015 Corporate Governance Code---are instruments of persuasion: they articulate norms for institutional investor engagement and board structure but impose no sanctions for non-compliance. Under voluntary frameworks, firms with the greatest incentives to resist reform---those in which controlling shareholders extract private benefits through related-party transactions---are precisely the ones least likely to opt in. The empirical record from Korea's own institutional-activism literature corroborates this: markets do not react to NPS vote-no announcements in the short run \citep{kim_nps_governance_firm_value}, suggesting that voluntary shareholder activism lacks the enforcement credibility necessary to shift valuation expectations. Japan's two early soft-code events therefore did not generate a sustained TOPIX re-rating that Korea failed to match, because neither code constituted an enforceable governance shock for the firms most responsible for Japan's pre-reform discount.

The 2023 TSE mandate crossed a qualitative threshold. By requiring firms with P/B below 1.0 to publicly disclose capital efficiency improvement plans or explain non-compliance---backed by the threat of delisting guidance and active TSE follow-up---the mandate transformed governance reform from a reputational norm into a compliance obligation with observable consequences. The immediate negative response in the MSCI EM Asia--TOPIX spread reflects investors anticipating differential capital allocation improvement: Japan's listed universe would improve while the broader EM Asia region, Korea included, would not. Korea's 2024 Value-Up program generated no comparable improvement: against the primary MSCI EM Asia benchmark, all three milestones produced persistently negative CARs of $-2.18$, $-1.80$, and $-2.12$ P/B points at the 12-month horizon---Korea was losing ground to regional peers who had themselves failed to keep pace with Japan.

The contrast between Japan's hard mandate and Korea's voluntary program is therefore not merely a question of intensity but of mechanism. Voluntary programs do not resolve the fundamental collective-action problem: controlling shareholders who benefit from the status quo have no individually rational reason to comply, and markets will not price in a reform they do not expect to be enforced. Whether the divergence---$-2.58$ P/B points in the primary MSCI EM Asia specification, $-6.48$ in the secondary KOSPI--TOPIX specification---fully closes the structural gap in the long run will depend on Japan maintaining enforcement and on Korea's eventual transition to comparable mandates.

\subsection{The Governance Channel and Its Limits}

The results of this paper are consistent with, but do not definitively establish, a governance channel as the exclusive mechanism behind the Korea Discount. Three alternative interpretations merit explicit consideration.

First, the discount could partly reflect structural differences in the industrial composition of KOSPI relative to TOPIX and MSCI EM Asia. Korean equities are heavily weighted toward manufacturing conglomerates with high tangible-asset intensity, precisely the firm characteristics identified as correlates of low P/B in cross-sectional regressions \citep{yang_korea_discount_growth_vs_governance}. Our index-level analysis cannot control for compositional differences of this kind, and the sub-period deterioration in Korea's mean P/B from 1.36 (pre-reform) to 0.93 (post-TSE) does not by itself isolate governance from slower structural transitions in Korean industry. However, the event-study design partially addresses this concern: if compositional differences are approximately stable within the 36-month pre-event window used for trend estimation, then the abnormal spread in the post-event period captures reform-related deviations rather than secular composition shifts. The consistency of negative CARs across three distinct Korea Value-Up milestones, spanning February through August 2024, is harder to explain through slow-moving compositional change than through a market-level assessment of Korea's governance trajectory.

Second, foreign exchange dynamics could conflate valuation and currency effects if P/B ratios are being compared across currency boundaries in a way that MSCI index methodology does not fully equalize. The panel OLS results partially address this concern: the log FX coefficient in the two-way fixed effects specification is $-0.024$ ($p = 0.018$), indicating a modest negative association between won depreciation and Korean P/B after controlling for country and time fixed effects, but the reform interaction terms for the TSE Corporate Governance Code and the TSE P/B Reform Request retain significance ($-0.322$, $p < 0.001$ and $-0.234$, $p = 0.056$ respectively), and the event-study design absorbs FX trends through the pre-event linear extrapolation. Currency effects alone cannot account for the persistence and magnitude of the divergence documented here.

Third, the growth-versus-governance debate in the literature---cross-sectional P/B regressions find that R\&D investment and firm maturity outperform governance variables \citep{yang_korea_discount_growth_vs_governance, lee_korea_discount_pbr}---raises the question of whether governance reform can raise valuations without also addressing the underlying growth deficits of large Korean conglomerates. Our results do not directly engage this question at the firm level, since we use index-level data. What the event-study evidence does establish is that Japan's market reacted positively to a governance mandate even though Japanese conglomerates share many of the same structural growth limitations as their Korean counterparts. If governance improvements were insufficient to raise valuations in the absence of fundamental growth, the TSE mandate should not have produced a sustained TOPIX re-rating. The magnitude of the 2023 CAR is therefore at least partially inconsistent with a pure growth story, and more consistent with the view that governance-induced capital allocation improvements---including buybacks, dividend increases, and divestiture of low-return subsidiaries---can drive re-rating independently of underlying business model transformation.

\subsection{Limitations}

Several limitations of the empirical design should be acknowledged. First, the stacked event-study methodology produces descriptive CARs rather than inferential estimates: the saturated cohort-by-relative-time specification places one observation per cell, precluding the construction of standard errors, and the clustering of Korea's three Value-Up milestones within a seven-month window creates calendar-time overlap that violates the standard stacked-DiD independence assumption. The CARs reported in Section~4 are therefore properly interpreted as measures of realized abnormal spread movement, not as causal effect estimates with associated uncertainty. Readers should be cautious about reading precision into the specific numerical values; the qualitative patterns---persistent negation in the MSCI EM Asia comparison, the sharp sign reversal between Japan's soft and hard codes---are robust to the design's descriptive limitations.

Second, Japan serves as a governance reform benchmark rather than a clean causal counterfactual. The identifying assumption underlying the primary Japan event study is that, absent Japan's governance reforms, the MSCI EM Asia--TOPIX spread would have continued along its pre-event linear trend. This assumption is implausible for long horizons---the indices respond to different macroeconomic shocks, exchange rate regimes, and sector compositions---but defensible within the 36-month estimation window used here. The secondary KOSPI--TOPIX analysis carries the same assumption for Korea specifically; the convergence of the primary EM Asia and secondary KOSPI results across all three reform events strengthens the inference that Japan's hard mandate, not Korea-idiosyncratic dynamics, drove the divergence.

Third, the analysis is conducted entirely at the index level and cannot speak to the cross-sectional distribution of governance improvements within Korea. The market response to Korea's 2025 Commercial Act amendment was heterogeneous: firms with the highest book-to-market ratios and the lowest controlling-shareholder ownership concentrations showed the largest abnormal returns, consistent with governance reform having its greatest effect precisely where expropriation risk is highest \citep{kang_commercial_act_market_reactions}. Our aggregate index-level design would not detect such heterogeneity even if it were present in the data, and the apparent absence of aggregate improvement in the KOSPI--MSCI EM Asia spread does not preclude meaningful governance gains at the firm level for a subset of Korean companies. Firm-level event studies around Korea's 2024 and 2025 reform milestones, using within-firm variation in governance exposure as a conditioning variable, would be a natural complement to the aggregate evidence presented here.

Fourth, the study window ends in April 2026, and only two post-event months are available for Korea's most recent mandatory reform events. The spaced follow-through specification suggests tentatively positive CARs for the July 2025 mandatory governance disclosure expansion and the February 2026 Value-Up dividend disclosure rule, but a 2-month window is entirely insufficient to distinguish a sustained structural improvement from short-run announcement effects. The conclusions of this paper pertain to the voluntary 2024 Value-Up phase; whether Korea's subsequent mandatory increments---and the broader 2025 Commercial Act fiduciary duty expansion---produce durable re-rating is an empirical question that cannot yet be answered with the available data.

\subsection{Policy Implications}

The evidence in this paper, read alongside the literature surveyed in Section~2, supports a specific policy conclusion: resolving the Korea Discount requires hard legislative reform that removes the voluntary character of governance compliance, rather than iterative enhancement of existing soft-law frameworks. The mechanism through which hard mandates operate is not primarily informational---markets are not learning new facts about Korean corporate governance from a mandatory disclosure rule that they could not have inferred from voluntary disclosures. Rather, mandatory frameworks change the strategic environment for controlling shareholders by making non-compliance costly and observable, thereby credibly committing the regime to enforcement in a way that voluntary guidelines cannot.

This distinction has direct relevance for Korea's reform agenda. The 2024 Corporate Value-Up program failed not because its objectives were poorly specified but because the voluntary mechanism left the participation decision with the very firms whose controlling shareholders had the largest private incentives to resist. A hard mandate analogous to the TSE P/B directive---requiring firms below book value to disclose capital efficiency improvement plans under threat of regulatory consequence---would remove that opt-out. The 2025 Korean Commercial Act amendment, which expanded directors' fiduciary duties from the corporation to shareholders broadly, represents meaningful progress: legislative passage produced significant positive abnormal returns for high-book-to-market firms, consistent with markets revising expropriation risk downward \citep{kang_commercial_act_market_reactions}. However, a fiduciary duty expansion creates a legal cause of action rather than a proactive compliance obligation; whether it shifts controlling shareholder behavior without accompanying exchange-level enforcement remains to be seen.

Governance reform and market openness are complementary instruments \citep{joo_foreign_investors_value_up}: foreign institutional investors impose discipline on both over-investing and under-returning firms, correcting the allocation distortions that flow from concentrated ownership structures. Reforms that reduce the Korea Discount through improved governance would therefore be self-reinforcing if they simultaneously attract foreign institutional participation, because greater foreign ownership raises the share of investors who monitor management aggressively and vote against related-party transactions. The Korean government's parallel efforts to improve market accessibility---including English-language corporate disclosures and simplified foreign investor registration---should be understood in this context as part of a governance reform package rather than a standalone capital-account liberalization measure.

The findings also bear on the broader debate over the relative contribution of governance failures and growth deficits to the Korea Discount. The evidence from Japan's 2023 mandate suggests that governance reforms can produce large, sustained valuation improvements even within an industrial structure broadly similar to Korea's---without requiring a transformation of the underlying business models of the large conglomerates that dominate both markets. This does not imply that growth-oriented reforms are irrelevant, but it does establish a lower bound on the valuation impact that governance reform alone can achieve. Policies that simultaneously address governance---through mandatory disclosure obligations, enhanced fiduciary duties, and exchange-level enforcement---and growth---through R\&D tax incentives, cross-shareholding reduction, and holding company discount compression---are likely to be complementary and mutually reinforcing, as the literature on Korea's chaebol structure suggests that governance improvements reduce value-destroying pyramidal acquisitions that otherwise crowd out productive investment.

\subsection{Conclusion}

This paper has argued and empirically documented that the Korea Discount is structural, governance-driven, and unlikely to close without hard legislative reform. The descriptive panel establishes KOSPI's mean discount of $-0.177$ P/B points to TOPIX and $-0.601$ to MSCI EM over 20 years, with a sub-period deterioration from near-parity in 2004--2013 to a 0.43 P/B point gap in 2023--2024. The geopolitical risk falsification rules out the North Korean threat explanation ($t = -0.835$, $p = 0.405$). The Japan event study identifies the 2023 TSE hard mandate as the structural inflection point: the primary MSCI EM Asia--TOPIX CAR reached $-2.58$ P/B points at $\tau = 20$, and the secondary KOSPI--TOPIX CAR reached $-6.48$ at $\tau = 24$---the largest event in the sample---while Japan's earlier soft codes produced positive CARs in both specifications. The Korea Value-Up event study shows the voluntary program produced no analogous improvement: the primary MSCI EM Asia specification records persistently negative CARs of $-2.18$, $-1.80$, and $-2.12$ P/B points at $\tau = 12$ for all three milestones, robust across event-date specifications and post-event window lengths.

The practical implication is that Korea currently stands where Japan stood before 2023: in possession of well-articulated soft-law governance norms whose voluntary character renders them structurally insufficient to overcome the incentive problems of concentrated ownership. Japan's experience demonstrates that the transition to hard exchange-level mandates can produce large, rapid, and sustained valuation improvements---improvements that were visible not just in Korea's relative performance but across the entire EM Asia regional basket. The 2025 Commercial Act amendment expanded the legal framework for shareholder protection in a way that markets have already partially priced, but an exchange-level enforcement mechanism---requiring firms below book value to disclose credible capital efficiency plans under threat of regulatory consequence---remains the missing complement. Until that transition occurs, Korea's valuation gap will persist as a structural equilibrium outcome of an ownership regime that voluntary guidelines have repeatedly failed to reform.
